Four syntactic analyses support the well-formedness rules: $\assignsS(t)$, the set of variables declared or assigned
anywhere in $t$ (without descending into nested function definitions, which have their own scope, so that
$\assignsS(t) = \dom(\Delta)$ when $t : \Assigns{\Delta}$); $\declaresS(t)$, the sequence of names declared in $t$, with
repeats, so that a scope can require its declarations to be distinct; $\binds(p)$, the set of variables introduced by a
pattern $p$, and likewise $\binds(\seq{g})$ by the generators of a qualifier sequence and $\binds(s)$ by the patterns within a statement; and $\captures(e)$, the
set of variables from the enclosing scope that are captured by lambdas within expression $e$, and likewise
$\captures(\seq{g})$ within a qualifier sequence.
All are defined in \secref{auxiliary-syntactic}.
